\chapter{Case Study: Building and Optimizing a Functional-Language Implementation}
\label{ch:functional-case-study}

The case study connects parsing, inference, optimization, closures, list cases, and execution in one small language.
It includes a strict source interpreter, a compiled stack machine, and an alternate lazy evaluator.
Independent execution paths make it possible to catch compiler defects rather than testing generated code only against itself.
The implementation is \code{example-code/functional-labs/fun.py} and uses only Python's standard library.

\section{Define the deliverable before extending it}
\label{sec:ch93-contract}
The language implements unary higher-order functions, let polymorphism, restricted recursive functions, integers, booleans, conditionals, and homogeneous lists.
It parses one validated S-expression and rejects trailing expressions.
The parser limits token count to 10,000 and nesting depth to 100; those are resource policies, not language-theoretic bounds.
Execution has a separate fuel counter and can also encounter host recursion and memory limits.

Table~\ref{tab:ch93-parts} maps the executed components.
The compiler emits structured stack code, not native machine code or an object file.
Its branch and binding instructions contain nested code blocks; the VM executes those blocks recursively on the host stack.
Consequently it does not promise proper tail calls or bounded host stack usage for source recursion.

\begin{table}[htbp]
\caption{The case study supplies several independent semantic paths but keeps its backend scope explicit.}
\label{tab:ch93-parts}
\centering
\begin{tabular}{ll}
\toprule
Component & Implemented responsibility \\
\midrule
parser/validator & complete S-expression, arities, binder restrictions \\
type solver & finite monotypes, schemes, occurs check, recursive placeholder \\
strict evaluator & lexical closures and left-to-right evaluation \\
optimizer & constant folding, selected identities, constant branch selection \\
compiler/VM & structured stack code, closures, recursion, list cases \\
lazy evaluator & thunks, memoization, black-hole detection, lazy fields \\
\bottomrule
\end{tabular}
\end{table}

The type checker runs before optimization in the CLI.
An ill-typed unreachable branch is therefore still rejected rather than optimized away before checking.
This gives the language a stable static contract independent of optimizer strength.
Tests also re-infer optimized expressions so a transformation that corrupts typing is caught separately from a runtime mismatch.

\section{Compile expressions to stack effects}
\label{sec:ch93-code}
Compiling a literal pushes one value; compiling a variable loads its bound cell's value.
Compiling add first emits the left operand, then the right, then ADD, which pops two values and pushes one.
Application follows the same operand order and APP enters a closure using its captured environment.
Every complete expression block must finish with exactly one stack value relative to its empty local operand stack.

Algorithm~\ref{alg:ch93-compile} describes the core compiler contract.
It assumes validated typed syntax and separates code for unselected branches into nested blocks.
Constructor cases likewise evaluate only the chosen branch at runtime.
The VM checks its final local stack shape rather than silently accepting extra or missing results.

\begin{algorithm}[htbp]
\caption{Compile a typed expression into structured code producing one value.}
\label{alg:ch93-compile}
\begin{pseudocode}
\Require validated typed expression and lexical binding contract
\Ensure code block with net one result under the strict semantics
\State literals and variables emit CONST or LOAD
\State lambdas emit CLOSURE with parameter and compiled body block
\State binary primitives emit left block, right block, then consuming opcode
\State applications emit function block, argument block, then APP
\State conditionals emit test block and IF holding two separate branch blocks
\State lets and cases retain their lexical binders and nested compiled blocks
\State verify stack effects and branch result compatibility
\end{pseudocode}
\end{algorithm}

An efficient code builder can emit $O(N)$ instructions for $N$ syntax nodes in this structured representation.
The reference compiler concatenates Python lists recursively, so a highly skewed expression can incur quadratic copying during code construction.
The chapter does not advertise the linear builder bound as a property of that implementation.
Replacing concatenation with an append-only builder is a straightforward extension whose emitted-code and evaluation-order tests should remain unchanged.

Figure~\ref{fig:ch93-stack} works add(4,7).
The stack evolves from empty to [4], then [4,7], then [11].
An incorrect compiler that evaluates the right side first would not be exposed by this commutative example alone.
Closure calls, conditionals, and termination-sensitive tests provide additional ordering evidence.

\begin{figure}[htbp]
\centering
\begin{tikzpicture}[node distance=13mm]
\node[basic block] (a) {CONST 4: [4]};
\node[basic block,below=of a] (b) {CONST 7: [4,7]};
\node[basic block,below=of b] (c) {ADD: [11]};
\draw[cfg edge] (a)--(b);\draw[cfg edge] (b)--(c);
\end{tikzpicture}
\caption{Primitive compilation preserves operand order and leaves one result on the stack.}
\label{fig:ch93-stack}
\end{figure}

\section{Implement a small sound optimizer}
\label{sec:ch93-opt}
The optimizer folds integer add, multiply, and integer equality when both operands are literal integers.
It removes add-zero and multiply-one wrappers and selects a branch when the test is a Boolean literal.
It traverses lambda bodies, let values and bodies, and list-case branches without confusing binders with expressions.
It does not inline, substitute by spelling, reassociate arithmetic, or erase multiplication by zero.

Algorithm~\ref{alg:ch93-opt} gives the safe local decisions.
Python's Boolean subtype relationship is again handled explicitly: a Boolean literal is not an integer constant for this source language.
Selected identities preserve evaluation of their nonconstant operand, while selecting a constant branch removes only a branch that the source would not evaluate.
These rules are sound under both the strict and lazy pure core contracts, though the two machines can differ on other programs.

\begin{algorithm}[htbp]
\caption{Apply local pure-core rewrites without erasing demanded divergent computation.}
\label{alg:ch93-opt}
\begin{pseudocode}
\Require typed pure-core expression with exact integer and Boolean literal kinds
\Ensure expression with the same permitted value/termination behavior
\State recursively optimize expression children under binding-aware syntax cases
\If{integer primitive operands are integer literals} \State fold exactly \EndIf
\If{add has right literal zero} \State replace by left operand \EndIf
\If{multiply has right literal one} \State replace by left operand \EndIf
\If{conditional test is a Boolean literal} \State retain only its selected arm \EndIf
\State otherwise rebuild the original operator with optimized children
\end{pseudocode}
\end{algorithm}

The traversal visits every source node once, excluding large-integer arithmetic and tuple/list construction costs.
Constant folding itself can increase integer bit length substantially, so a production compiler needs a constant-size policy as well as an expression-node budget.
The small lab's input and evaluation limits do not constitute a complete denial-of-service defense for hostile numeric inputs.
The README records that boundary rather than calling the implementation a hardened service.

\section{Work recursion and constructor matching}
\label{sec:ch93-examples}
The checked-in sum program recursively matches a list and adds its head to the sum of its tail.
On [4,7,9], the recursive results are zero at Nil, nine at the final Cons, sixteen at the preceding Cons, and twenty at the root.
Table~\ref{tab:ch93-sum} records the complete return sequence.
Strict interpretation, compiled VM execution, and lazy scalar evaluation all return twenty on this terminating pure input.

\begin{table}[htbp]
\caption{The list-sum case study exercises recursion, constructor dispatch, and field binding.}
\label{tab:ch93-sum}
\centering
\begin{tabular}{lrr}
\toprule
Remaining list & Head contribution & Returned sum \\
\midrule
empty & 0 & 0 \\
$[9]$ & 9 & 9 \\
$[7,9]$ & 7 & 16 \\
$[4,7,9]$ & 4 & 20 \\
\bottomrule
\end{tabular}
\end{table}

The factorial example applies a recursively defined function to six and returns 720.
The base test demands an integer equality; the recursive branch decrements by adding negative one.
The recursion is not tail-recursive because multiplication remains after the recursive call.
The language has no special subtraction syntax and no hidden native factorial primitive.

The lexical capture test returns fifteen despite a later x=99 binding.
The polymorphic identity test applies one let-bound function at Bool and Int types.
The lazy sharing test counts one first-time force for x=18+5 even though x appears twice.
Each example targets a different mechanism that a list-sum output alone would not distinguish.

\section{Compare independent execution paths}
\label{sec:ch93-tests}
For every a and b from negative six through six and each of add, mul, and eq, the suite constructs a closed higher-order program calling a lambda on a.
That yields 507 generated programs.
It compares strict source execution, compiled VM execution, optimized compiled execution, lazy execution, and optimized lazy execution against independently computed host arithmetic expectations.
Type inference runs before execution and after optimization.

Figure~\ref{fig:ch93-oracles} separates source and compiled paths.
The tests also exercise deliberate type errors, malformed syntax, recursive restrictions, original examples, CLI modes, and the strict/lazy termination difference.
They do not infer semantic equivalence from two fuel-exhausted runs.
The first-order rewrite lab separately checks trace replay and graph hazards rather than being folded into one opaque end-to-end test.

\begin{figure}[htbp]
\centering
\begin{tikzpicture}[node distance=13mm and 13mm]
\node[basic block] (source) {typed source expression};
\node[basic block,below left=of source] (eval) {strict source interpreter};
\node[basic block,below right=of source] (vm) {compiled structured VM};
\node[basic block,below=of source,yshift=-20mm] (compare) {compare result and independent expectation};
\draw[cfg edge] (source)--(eval);\draw[cfg edge] (source)--(vm);
\draw[cfg edge] (eval)--(compare);\draw[cfg edge] (vm)--(compare);
\end{tikzpicture}
\caption{Independent execution paths expose defects that self-comparison can miss.}
\label{fig:ch93-oracles}
\end{figure}

\section{Extend the implementation without broadening claims accidentally}
\label{sec:ch93-next}
General algebraic types require declarations, constructor schemes, layout metadata, and pattern coverage.
Native compilation requires concrete closure layouts, calling conventions, allocation, stack maps, and code emission.
Proper tail calls require an explicit continuation or frame-reuse mechanism rather than recursive calls into Python.
A production lazy runtime additionally needs concurrent thunk policy, exception memoization semantics, profiling, and a collector designed for its heap objects.

These extensions exceed the tested core; the lab is not a production compiler or optimal graph reducer.

\section{State the compiler simulation invariant}
\label{sec:ch93-simulation}
For the strict fragment, the central compiler invariant relates source environments to VM environments containing corresponding binding cells.
Compiling an expression and executing its code should produce one stack result corresponding to the source value, while preserving the lexical meaning of every captured binding.
For constants and variables this follows directly from CONST and LOAD.
For primitives, the induction hypotheses for the two operand blocks supply their values in source order before the consuming instruction runs.

Lambda compilation needs the relation to extend to closures: source syntax bodies correspond to compiled body blocks, and their captured environments must relate field by field or binding by binding.
Application then uses those related closure environments, not the caller's unrelated shadowing bindings.
Recursive let adds a pair of cells to the relation before constructing the closures, then fills both with their corresponding closure values.
This explains why the recursive lambda restriction is part of the implementation proof rather than just a parsing convenience.

Conditionals and cases select one related branch block producing one result.
Stack assertions and differential tests check instances and structural consequences of this invariant, not a mechanized proof of the full simulation.

\section{Exercises}
\label{sec:ch93-exercises}
\begin{enumerate}
\item\label{ex:ch93-code} Compile add(4,7), showing instructions and stack after each instruction.
\item\label{ex:ch93-sum} Run the list-sum program on [2,5], giving every recursive return value.
\item\label{ex:ch93-opt} Explain why folding mul(app(loop,0),0) to zero is intentionally forbidden in the strict optimizer.
\item\label{ex:ch93-backend} Name the runtime and backend components needed to replace the structured Python VM with native code while retaining closures and recursion.
\end{enumerate}

\section{Further reading}
\label{sec:ch93-reading}
Appel supplies a larger functional compiler architecture~\cite{appel-1998-modern}.
Peyton Jones develops a compiled lazy runtime in detail~\cite{peytonjones-1992-stg}, and Shivers explains optimization-oriented higher-order control-flow recovery~\cite{shivers-1991-control}.
The lab README and exercise solutions document reproducible execution and the remaining implementation boundaries.
